\phantomsection
\color{unimain}\section*{\centering{\large{ПЕРЕЧЕНЬ СОКРАЩЕНИЙ}}}
\addcontentsline{toc}{section}{Перечень сокращений}
{\color{white}\fontsize{0.1pt}{0.1pt}\selectfont<begABBRS>}
\color{black}
\begin{longtable}{>{\raggedright\arraybackslash}m{2cm}>{\raggedright\arraybackslash}m{0.5cm}>{\raggedright\arraybackslash}m{20cm}}
\endfirsthead\endhead\endfoot\endlastfoot
GOOSE & -- & Generic Object-Oriented Substation Event; \\
IRF & -- & Internal Relay Fault; \\
PPS & -- & Pulse Per Second; \\
PTP & -- & Precision Time Protocol (протокол точного времени); \\
SNTP & -- & Simple Network Time Protocol (простой сетевой протокол времени); \\
АПВ & -- & автоматическое повторное включение; \\
АСУ~ТП & -- & автоматизированная система управления технологическими процессами; \\
АУВ & -- & автоматика управления выключателем; \\
АЦП & -- & аналого-цифровой преобразователь; \\
БВКЗ & -- & блокировка при внешних коротких замыканиях; \\
БДО & -- & быстродействующий дифференциальный орган; \\
БИ & -- & блок испытательный; \\
БКУ & -- & блок команд управления; \\
БНН & -- & блокировка при неисправности цепей напряжения; \\
БНТ & -- & блокировка при бросках намагничивающего тока; \\
БУ & -- & блок управления; \\
ВН & -- & высшее напряжение; \\
ВЧ & -- & высокочастотный; \\
ВЧПП & -- & высокочастотный приемо-передатчик; \\
ДЗ & -- & дистанционная защита; \\
ДЗО & -- & дифференциальная защита ошиновки; \\
ДЗШ & -- & дифференциальная защита шин; \\
ДО & -- & дочерние общества; \\
ДОТХ & -- & дифференциальный орган с торможением; \\
ДТО & -- & дифференциальная токовая отсечка; \\
ДФЗ & -- & дифференциально-фазная защита; \\
ЕЭС & -- & единая энергетическая система; \\
ЗАПВ & -- & запрет автоматического повторного включения; \\
ЗИП & -- & запасные части, инструменты и принадлежности; \\
ЗНР & -- & защита от неполнофазного режима; \\
ЗНФ & -- & защита от непереключения фаз; \\
ИО & -- & измерительный орган / избирательный орган; \\
ИЧМ & -- & интерфейс человек-машина; \\
КЗ & -- & короткое замыкание; \\
КОН & -- & контроль отсутствия напряжения; \\
КС & -- & контроль синхронизма / контрольная сумма; \\
КСВ & -- & контроль силового выключателя; \\
КСН & -- & контроль синхронизма и напряжения; \\
КЦН & -- & контроль цепей напряжения; \\
КЦТ & -- & контроль цепей тока; \\
ЛО & -- & логика отключения; \\
ЛОП & -- & логика опробывания присоединений; \\
ЛОЧ & -- & логика очувствления; \\
ЛЭП & -- & линия электропередачи; \\
МДО & -- & медленнодействующий дифференциальный орган; \\
МЭК & -- & международная электротехническая комиссия; \\
НВЧЗ & -- & направленная высокочастотная защита; \\
НН & -- & низшее напряжение; \\
ННЭ & -- & наличие напряжения на энергообъекте; \\
ОАПВ & -- & однофазное автоматическое повторное включение; \\
ОЗУ & -- & оперативное запоминающее устройство; \\
ОНЭ & -- & отсутствие напряжения на энергообъекте; \\
ОО & -- & орган опробования / общая обмотка; \\
ООЧ & -- & орган очувствления; \\
ОП & -- & обратная последовательность; \\
ОТ & -- & оперативный ток; \\
ПА & -- & противоаварийная автоматика; \\
ПЗУ & -- & постоянное запоминающее устройство; \\
ПК & -- & предохранительный клапан; \\
ПО & -- & пусковой орган / программное обеспечение; \\
ПРД & -- & передатчик; \\
ПРМ & -- & приемник; \\
ПУ & -- & пульт управления / поперечное ускорение; \\
РД & -- & реле давления; \\
РЗА & -- & релейная защита и автоматика; \\
РКВ & -- & реле команды «включить»; \\
РКО & -- & реле команды «отключить»; \\
РС & -- & регистрация событий / реле сопротивления; \\
РТ & -- & реле тока; \\
РУ & -- & распределительное устройство; \\
РФ & -- & реле фиксации; \\
РФК & -- & реле фиксации команд; \\
РФП & -- & реле фиксации включенного положения; \\
РЭ & -- & руководство по эксплуатации; \\
СО & -- & система охлаждения; \\
ССПИ & -- & система сбора и передачи информации; \\
СШ & -- & система (секция) шин; \\
ТЗНП & -- & токовая защита нулевой последовательности; \\
ТН & -- & трансформатор напряжения; \\
ТО & -- & токовая отсечка; \\
ТР & -- & технический регламент; \\
ТС & -- & технологическая сигнализация; \\
ТТ & -- & трансформатор тока; \\
ТУ & -- & телеуправление / телеускорение / технические условия; \\
УВ & -- & управление выключателем / управляющее воздействие; \\
УРОВ & -- & устройство резервирования при отказе выключателя; \\
УС & -- & улавливание синхронизма; \\
ФБ & -- & функциональный блок; \\
ФПВ & -- & фиксация положения выключателя; \\
ФСУ & -- & функциональная схема устройства; \\
ЦВ & -- & цепи включения; \\
ЦН & -- & цепи напряжения; \\
ЦО & -- & цепи отключения; \\
ЦОС & -- & цифровая обработка сигналов; \\
ЦП & -- & центральный процессор; \\
ЦПС & -- & цифровая подстанция; \\
ЦТ & -- & цепи тока; \\
ЦУ & -- & цепи управления; \\
ЧТО & -- & чувствительный токовый орган; \\
ШОН & -- & шкаф отбора напряжения; \\
ШП & -- & шинки питания; \\
ШЭТ & -- & шкаф электротехнический типовой; \\
ЭМВ & -- & электромагнит включения; \\
ЭМО & -- & электромагнит отключения; \\
ЭМУ & -- & электромагнит управления. \\
{\color{white}\fontsize{0.1pt}{0.1pt}\selectfont<endABBRS>}
\end{longtable}
